\subsection{Numerical Integration Schemes}\label{subsec:num_integration}

This section reviews the primary numerical integration methods used for computing observables.

\subsubsection{Quadrature-based integration}

For 1D integrals of the form:

\begin{equation}
I = \int_a^b f(x) dx
\end{equation}

approximate as weighted sum of function values:

\begin{equation}
I \approx \sum_{i=1}^n w_i f(x_i)
\end{equation}

where $x_i$ are quadrature nodes and $w_i$ are weights.

\subsubsection{Common quadrature rules}

\begin{description}
    \item[Gauss-Legendre]: Optimal for smooth functions; weights/nodes determined by orthogonal polynomials
    \item[Gauss-Laguerre]: For integrals $\int_0^\infty e^{-x} f(x) dx$ (natural exponential decay)
    \item[Adaptive (QuadGK)]: Subdivides domain; refines where error is large
    \item[Clenshaw-Curtis]: Uses Chebyshev nodes; fast FFT-based evaluation
\end{description}

\subsubsection{Composite integration}

For integrals with multiple peaks or wide dynamic range:

\begin{equation}
I = \sum_k \int_{a_k}^{b_k} f(x) dx \approx \sum_k Q_k[f]
\end{equation}

Apply quadrature to each sub-domain $[a_k, b_k]$, then sum.

\subsubsection{Cosmological examples}

Many cosmological integrals have special structure:

\begin{itemize}
    \item \textbf{Comoving distance}: $\chi(z) = \int_0^z \frac{dz'}{E(z')}$ → smooth, monotonic integrand
    \item \textbf{Lensing kernels}: $W_\ell(k) = \int d\chi \, \phi(\chi) j_\ell(k\chi)$ → oscillatory Bessel function
    \item \textbf{Power spectrum convolution}: $\int dk k^2 P(k) W(k)$ → power-law decay
\end{itemize}

Each requires adapted strategy (see next sections).

See \cref{subsec:num_quadrature} for specific quadrature algorithms.
